%% =============================================================================
%% Section 8 — Conclusion
%% "How Do LLM-Based Agents Represent Code Repositories?"
%% ACM Transactions on Software Engineering and Methodology (TOSEM)
%% =============================================================================

\section{Conclusion}
\label{sec:conclusion}

This paper presented the first Systematic Literature Review dedicated to the
\emph{Repository Representation Problem} (RRP) — the challenge of selecting,
encoding, and injecting the right subset of a large codebase into a finite
LLM context window.
Analysing 159 primary studies (2020--2026) through a rigorous five-pillar
methodology (Kitchenham~\cite{kitchenham2007guidelines},
Wohlin~\cite{wohlin2014snowballing}, PRISMA~2020~\cite{page2021prisma},
Cruzes and Dyb{\aa}~\cite{cruzes2011recommended},
ACM SIGSOFT Empirical Standards~\cite{acmsigsoft2020standards}), we
produced six interlocking contributions:

\begin{enumerate}

  \item \textbf{A seven-paradigm taxonomy (PAR-1 through PAR-7)} grounded in
        the corpus, with formal definitions and unambiguous paradigm boundaries.
        The taxonomy spans dense code embeddings (PAR-1), retrieval-augmented
        context construction (PAR-2), graph-based structural representations
        (PAR-3), execution and dynamic representations (PAR-4), memory-augmented
        context (PAR-5), compression and summarization (PAR-6), and agentic
        iterative exploration (PAR-7).

  \item \textbf{A ten-dimension analytical framework} (context type,
        granularity, construction pipeline, complexity, token budget,
        freshness, semantic precision, scalability, cross-file dependency
        support, training dependency) that enables systematic paradigm-level
        comparison — a level of analytical resolution absent from all five
        existing adjacent surveys.

  \item \textbf{A 7$\times$7 paradigm--SE-objective mapping matrix} showing
        which paradigms serve which software engineering tasks (code completion,
        code generation, program repair, fault localisation, code search, code
        translation, test generation), with three actionable findings: PAR-7 is
        the only universal paradigm; PAR-6 is absent from program repair;
        PAR-1 is concentrated almost exclusively in code search.

  \item \textbf{Empirical evidence of a 10--61\% static-repository evaluation
        bias} affecting PAR-1, PAR-3, and PAR-6 systems, as demonstrated by
        HumanEvo~\cite{zheng2025humanevo}.
        This finding reframes repository representation as a
        \emph{temporal engineering problem}, not merely a context selection
        problem.

  \item \textbf{Five structural research gaps} with a prioritised research
        agenda: temporal freshness (closing the static-snapshot gap), token-
        budget-aware paradigm hybridisation, multilingual repository support,
        evaluation standardisation, and privacy-aware local representation.

  \item \textbf{A fully documented PRISMA~2020 protocol} (flow diagram,
        inclusion/exclusion log, IRR $\kappa = 0.81$) enabling future
        replication and incremental updates as the field evolves.

\end{enumerate}

\subsection{The Added Value Over Existing Surveys}
\label{subsec:addedvalue}

The five most closely related secondary studies — Hou et al.\
\cite{hou2024llmse_slr}, Husein et al.~\cite{husein2025llmcompletion_slr},
Jiang et al.~\cite{jiang2026surveycodegen}, Tao et al.\
\cite{tao2026ragsurvey}, and Zhang et al.~\cite{zhang2026surveyse} — each
represent landmark contributions to the LLM-for-SE literature.
This SLR does not compete with them; it is orthogonal to them.

Those surveys collectively answer the question: \emph{what can LLM-based
systems do for software engineering?}
This survey answers a complementary question: \emph{how do LLM-based systems
see the repository when they do it?}

The distinction is architecturally significant.
Every LLM-based SE system, regardless of its target task, must first solve the
RRP: it must decide how to represent the repository as context before it can
generate, repair, search, or translate code.
Treating this representational choice as an implementation detail — as the five
surveys above do by necessity — leaves practitioners without a principled
framework for design decisions and leaves researchers without a vocabulary for
cross-study comparison.

This SLR establishes that vocabulary.
It demonstrates that repository representation is not a solved sub-problem of
prompt engineering.
It is a first-class architectural problem with at least seven structurally
distinct solution families, non-trivial operational trade-offs, and an open
empirical question about evaluation validity that affects the entire field.

The five surveys provide the landscape of what LLM-based SE agents accomplish.
This SLR provides the \emph{map} of how they reach their destination.
Both are necessary for the next generation of repository-aware SE tools; used
together, they give practitioners and researchers a complete picture of the
LLM-for-SE design space.

\subsection{Research Agenda}
\label{subsec:agenda}

Based on the five structural gaps identified above, we prioritise three
research directions for immediate attention:

\paragraph{Priority~1: Temporal freshness and live-repository representations.}
Closing the 10--61\% gap between static-snapshot evaluation and real-world
deployment requires both new benchmark designs with explicit temporal controls
and new representation techniques that support incremental, low-latency updates
for PAR-1, PAR-3, and PAR-6 systems.
Future work should replicate the HumanEvo~\cite{zheng2025humanevo} methodology
across multiple programming languages and repository types to obtain a robust
cross-paradigm bias estimate.

\paragraph{Priority~2: Token-budget-aware paradigm hybridisation.}
The complementary strengths of PAR-2 (scalability), PAR-3 (structural
precision), PAR-5 (longitudinal memory), and PAR-7 (flexibility) suggest
that hybrid paradigm compositions will define the next generation of
repository-level LLM tools.
A formal framework for composing paradigms within a token budget $\tau$ —
analogous to multi-objective query planning in database systems — is the most
pressing theoretical contribution the community can make.

\paragraph{Priority~3: Evaluation standardisation across paradigms.}
The field currently lacks a shared evaluation framework that can fairly compare
systems across all seven paradigms.
We call for a community-driven benchmark effort designed to control for
temporal bias~\cite{zheng2025humanevo}, repository diversity
\cite{ding2023crosscodeeval}, and task coverage
\cite{jimenez2024swebench,liu2024repobench} simultaneously — establishing for
repository representation research the same evaluation rigour that SWE-bench
\cite{jimenez2024swebench} established for repository-level issue resolution.
